\chapter{Conclusion and open problems}\label{ch:conclusion}

This thesis asked whether safety in a market with adaptive agents requires knowing who the agents are. For the markets studied here, it does not. Prices reveal the agents only through a merged measure of their memories and weights, and recovering the agents from that measure costs trading capital of order $\delta^{-4}$, experiments of order $(\Delta\eta)^{-(4K-2)}$, and with general memory an accuracy that improves only logarithmically (Chapter~\ref{ch:p1}). Safety depends on the agents only through their total feedback weight, which is certified at the parametric rate under a memory bound and cannot be certified without one; a guard built on that certificate is the largest safe admission rule (Chapter~\ref{ch:p2}). The theory of Chapter~\ref{ch:p3} places these results in one picture, in which the cost of a question is set by its behaviour at long memory, and shows that the safety bound needs only the agents' total sensitivity, even when their rules are nonlinear and their memories and weights change, and that learning makes certificates perishable at a computable rate. In data, disclosed trader classes reveal what prices do not, and the safety statistic forecasts adverse excursions in thin books, while the mechanism behind that forecast remains unidentified.

The policy message is direct. Supervisors who want to control transient risk should ask for the information that certifies safety, a bound on memory and on total sensitivity, renewed as often as the agents learn, rather than for the identities of the agents. Identity-based regimes may serve other purposes, such as enforcement and attribution, but this thesis shows they are neither necessary nor efficient for safety.

\section{Evidence that would settle the thesis}\label{sec:evidence}

The empirical work of Chapters~\ref{ch:p1} and~\ref{ch:p2} tests each paper's own predictions. This section states what that evidence establishes for the thesis as a whole, what it does not, and two tests designed to test the thesis directly. Both are registered: each specification was pushed to the public repository in a commit of its own before any of its data were downloaded. Neither has been run.

\subsection{What the existing evidence establishes}\label{sec:ev_existing}

\paragraph{Prices do not reveal the agents.} In 30 Binance perpetuals, trader-class positions pin down memory rates 28 to 124 times more precisely than prices and flow, and in a pre-registered test the CFTC's trader-class reports do the same in all eight regulated futures markets examined (Sections~\ref{p1:sec:disc_data} and~\ref{p1:sec:cot}). This is the identification half of the thesis in data: the merged measure leaves class memories essentially unrestricted.

\paragraph{The safety statistic carries information.} In thin Binance books, six-hour estimates of the amplification margin forecast the size of the next hour's adverse excursion out of sample in every pre-registered forecasting test, beyond a full volatility-forecasting baseline and VPIN; the effect is about a fifth as large in thin US stocks, negligible on Bybit and absent on Coinbase (Section~\ref{p2:sec:thin} and Table~\ref{p2:tab:alltests}).

\paragraph{What it does not establish.} The forecasting result is reduced form. Tests of a leverage or liquidation mechanism fail (Section~\ref{p2:sec:p22}), and the follow-up that points to slow reversion is close to a check of how the statistic is built (Section~\ref{p2:sec:p25}). As a risk control, the certified guard improves on a static limit but not, at matched exposure, on a volatility cap (Section~\ref{p2:sec:p24}). And from passive data the certificate itself is not always precise: in Paper~1's registered test the summed response was significantly different from zero in only 5 of 30 thin markets (Section~\ref{p1:sec:disc_data}). None of the existing tests checks the thesis's central claim directly: that, on the same data, certifying is much cheaper than identifying.

\subsection{Test E1: certify versus identify on fresh data}\label{sec:ev_e1}

\paragraph{Idea.} The theory says that the summed response is a linear functional of the merged measure and is cheap to certify, while the agents' memories are expensive to identify. In passive market data both can be estimated from the same regression of order flow on lagged returns: the sum of its coefficients is the certificate quantity $\hat\Phi$, and two exponentials fitted to the same coefficients give the memory rates. Published trader-class positions give a third, class-tagged estimate of memory. The test compares the precision of the three on data not used before.

\paragraph{Design.} Thirty Binance USD-M perpetuals (those of Section~\ref{p1:sec:disc_data}) on 105 fresh dates: the 4th, 5th, 18th, 19th and 23rd of every month from January 2025 to September 2026. A scan of every result file, notebook and script in the repository finds no earlier use of these symbol-dates. From five-minute returns, signed taker flow and the three class ratios, with Paper~1's estimators and day-bootstrap intervals, each symbol gives $I_{\rm cert}$, the width of the interval for $\hat\Phi$ relative to $|\hat\Phi|$, and $I_{\rm id}$, the width of the interval for the separation $\hat r_{\rm slow}-\hat r_{\rm fast}$ of the price-only memory rates relative to the separation. Primary~1: $I_{\rm id}\ge10\,I_{\rm cert}$, that is, certifying is at least ten times cheaper than identifying, in at least 75\% of symbols. Primary~2: the price-only memory intervals are at least ten times wider than the class-tagged ones in at least 75\% of symbols, a replication of Paper~1 on new days. Each primary must also pass a one-sided binomial test, Holm-corrected across the two. Primary~1 is a risky prediction: Paper~1's earlier test suggests that $\hat\Phi$ may be imprecise in thin markets, in which case it fails, and the failure will be reported. It would not contradict the theorems, which concern controlled experiments, but it would limit the empirical claim to what passive data can deliver.

\paragraph{Status.} Registered on 24 October 2026 (specification hash \texttt{0b6c31c0\ldots}; pipeline, download code and a Colab notebook committed with it). The notebook downloads the data, opens the test once and writes the results, the numbers printed here and a log of the SHA-256 of every downloaded file.

\subsection{Test E2: a prospective test of the mechanism}\label{sec:ev_e2}

\paragraph{Idea.} Every pre-registered test so far used archived data. A prospective test removes any doubt about the order of registration and data: the data do not exist when the specification is written.

\paragraph{Design.} The original specification (pushed on 23 October 2026) asks, for 20 thin perpetuals over at least three months after registration, whether six-hour $\hat C_3$ forecasts next-hour liquidation volume beyond volatility, impact and lagged liquidations (Primary~1), and whether adding realized next-hour liquidations to the excursion model removes at least half of $\hat C_3$'s forecasting gain (Primary~2). Binance's forced-liquidation stream, which it named as the source, cannot be recorded from the United States, and the public archive contains no liquidation files. An addendum, pushed on its own before any post-registration file was downloaded, keeps the hypotheses and replaces liquidation volume by forced deleveraging measured from open interest: the open interest destroyed during five-minute moves larger than twice their recent standard deviation. It also fixes the window (24 October 2026 to 31 January 2027), the split into 50 training and 50 test dates, and the estimators, which are those of the thin-book tests of Section~\ref{p2:sec:thin}. The price of the substitution is stated in advance: E2 can no longer test the liquidation channel with direct liquidation data, which was its advantage over the open-interest proxy of Section~\ref{p2:sec:p22}.

\paragraph{Status.} The analysis code and its notebook are committed with the addendum and refuse to run before 2 February 2027.

\subsection{Withdrawn results}\label{sec:ev_withdrawn}

Two results reported in an earlier draft of this thesis are withdrawn. The first was presented as a run of Test~E1 in 38 markets. It did not estimate anything from fresh data: it read results archived from earlier tests and used hard-coded and invented intervals, and its design compared quantities in different units. The specification it ran under is superseded by the one above, and none of its numbers appears in this thesis. The second was a simulation of learning agents whose weights, by construction, could only increase; its ``realized lifetimes'' measured that construction and could not have fallen below the bound they were said to test. Section~\ref{sec:bl_ai} replaces it with a registered simulation in which weights follow the agents' trading profits. The code and output of both are kept in the repository, marked as withdrawn, so that the record is complete.

\subsection{Rules for every new test}\label{sec:ev_rules}

\begin{enumerate}
\item The specification, its hash and an unlock counter are pushed to the public repository in a commit of their own, so the platform records when they became public, before any test data are downloaded or recorded.
\item Each test has at most two primaries with a pre-stated pass rule and a pre-stated correction across them.
\item A failed primary is reported as failed. An explanation found afterwards is labelled post hoc, and any test that follows from it is labelled as registered after the failure.
\item Every number in the thesis is produced by code that writes a result file, and an automated check compares the printed numbers with those files.
\item A result found to be wrong after it was reported is withdrawn in the text, with the reason, and its files are kept and marked rather than deleted.
\end{enumerate}

\section{Open problems}

\begin{enumerate}
\item \emph{Nonlinear questions.} Theorem~\ref{thm:dich}(iv) gives the exact cost of every linear question. Extend it to questions that are not linear in the merged measure, beginning with the peak of the impulse response, and make the link to the $1/\log N$ rate of identification (Corollary~\ref{cor:cases}(e)) a theorem rather than a reading of Chapter~\ref{ch:p1}'s proof.
\item \emph{Constants.} Theorem~\ref{thm:philog} fixes the order $\sigma^2\epsilon^{-2}\log(1/\epsilon)$ of the cost of certifying safety but not its constant, which depends on a harmonic measure; compute it, or the limit of $\omega_L(\delta)/(\delta\sqrt{\log(1/\delta)})$, as a function of the memory bound.
\item \emph{Minimal disclosure for safety.} Characterize the least informative disclosure under which a given safety question becomes certifiable from finite data. Theorem~\ref{thm:dich}(iii) says a memory bound is necessary for the summed response; is it sufficient in a sense that can be made operational?
\item \emph{Strategic agents.} Remark~\ref{rem:strategic_counterexample} refutes existence under a discontinuous admission rule. Determine sufficient assumptions for existence under a regularized guard and for the claimed comparison of equilibrium sensitivity.
\item \emph{Observational certification.} Develop finite-sample certificates from observational order-book data with dependence and non-stationarity, replacing the bootstrap diagnostics of Chapter~\ref{ch:p2}; Test~E1 measures how far passive data are from the reset contract of Theorem~\ref{thm:dich}.
\item \emph{The mechanism of thin-book risk.} Run Tests E1 and E2 of Section~\ref{sec:evidence}.
\end{enumerate}
